% timeFont 論文誌原稿（未来日記方式）。
% 組み立ては ./build.sh draft（下書き版）、./build.sh plain（印を一切見せない版。tex の印は残る）、./build.sh final（投稿版）。
% 未来日記の印: \begin{fd}{id}{status}{owner}{desc} ... \end{fd} と \fdcaption{id}{status}{owner}{desc}{caption}。
% 印の付いた記述は、まだ事実になっていない。事実台帳は facts_ledger.md、数値の設計は future_diary_plan.md にある。
\documentclass[submit]{ipsj}
\usepackage[dvipdfmx]{graphicx}
\usepackage{url}
\usepackage{amsmath}

%%%%% 未来日記の仕組み %%%%%
\newif\ifFDfinal
\newif\ifFDplain
\providecommand{\figsdir}{figs}
\IfFileExists{fdmode.tex}{\input{fdmode}}{\FDfinalfalse\FDplainfalse}
\definecolor{fdblue}{rgb}{0.10,0.25,0.55}
\definecolor{fdred}{rgb}{0.75,0.10,0.10}
\newcommand{\fdtag}[2]{\textcolor{fdred}{\textbf{\footnotesize ［未来日記 #1・#2］}}}
\def\fdDone{done}
\newcommand{\fdcheck}[2]{\ifFDfinal\def\fdStatus{#2}\ifx\fdDone\fdStatus\else
  \ClassError{ipsj}{Future diary #1 is not done}{Resolve it or build in draft mode.}\fi\fi}
\newenvironment{fd}[4]{\fdcheck{#1}{#2}\ifFDfinal\else\ifFDplain\else\par\begingroup\color{fdblue}\fdtag{#1}{#2}\fi\fi}%
  {\ifFDfinal\else\ifFDplain\else\endgroup\par\fi\fi}
\newcommand{\fdcaption}[5]{\fdcheck{#1}{#2}\ifFDfinal\caption{#5}\else\ifFDplain\caption{#5}\else\caption{\fdtag{#1}{#2}#5}\fi\fi}
%%%%%%%%%%%%%%%%%%%%%%%%%%%%

\setcounter{巻数}{68}
\setcounter{号数}{1}
\setcounter{page}{1}

\begin{document}

\title{認識率の曲線に合わせて文字が現れる動的フォント timeFont による\\
聞く人と見る人の公平なかるた対戦：上毛かるたへの適用と評価}
\etitle{timeFont: Fair Karuta Matches between Listeners and Viewers\\
with a Dynamic Font Revealing Letters at the Rate of Speech Recognition}

\affiliate{TSUDA}{津田塾大学\\
Tsuda University, Kodaira, Tokyo 187--8577, Japan}

\author{栗原 一貴}{Kazutaka Kurihara}{TSUDA}[kurihara@tsuda.ac.jp]
\author{丸山 礼華}{Reika Maruyama}{TSUDA}

\begin{abstract}
かるたでは，読み手の声が先頭から少しずつ届き，聞く人は第一音のうちに札を絞り込む．
文字を見て札を当てる人が同じ場で競うには，文字が見えてくる速さを，声が分かっていく速さに合わせて設計する必要がある．
本稿では，音声を聞いて音が分かっていく速さに合わせて文字が見えてくる速さを設計した動的フォント timeFont を用い，
聞く人と見る人が公平に早押しを競えるかるたの対戦を構成し，群馬県の郷土かるたである上毛かるた（44 枚，すべて一字決まり）で評価する．
音声と文字のそれぞれについて，提示を途中で打ち切ったときの認識率の曲線を実測し，ロジスティック曲線のレベル $x_0$ と傾き $\beta_1$ で要約する．
公平にする仕組みとして，速さを 5 段階に変えて実測したフェードの候補から札ごとに $\beta_1$ の最も近い段階を選択する方法（候補版），
候補版から内挿した提示を測定と更新の反復で音声の曲線に収束させる方法（精密版），
および押した時刻の早押し度 $1-f(t)$ で勝敗を判定するアプリの 3 つを構成した．
クラウドソーシングによる測定から，候補版では音声との認識率の差の最大が全札で 0.046 以内，精密版では反復 2 回で 0.034 以内であった．
聴者 20 名（10 組）による対戦実験では，補正のない提示で聞く役が 20 試合中 19 試合に勝ったのに対し，精密版の提示とアプリでは取り札数の差が消え，公平だと感じたかの評定も中立に近い値になった．
\end{abstract}

\begin{jkeyword}
動的フォント，かるた，聴覚障害，アクセシビリティ，認識率，心理測定関数，公平性
\end{jkeyword}

\begin{eabstract}
In karuta, a player who listens to the reader and a player who reads text cannot compete on equal terms, because speech reaches the listener incrementally whereas text appears at once or at a pace unrelated to speech.
This paper presents a fair match between listeners and viewers using timeFont, a dynamic font whose letters become legible at the same rate as speech becomes recognizable, and evaluates it on Jomo Karuta, a regional card game with 44 cards, each determined by its first syllable.
We measure recognition-rate curves for truncated speech and for partially revealed letters, summarize each curve by the level $x_0$ and the slope $\beta_1$ of a logistic fit, and construct three fairness mechanisms: a candidate version that selects the fade speed closest in $\beta_1$ from five measured speeds, a refined version that converges to the speech curve through measure-and-update iterations, and an application that judges by the earliness score $1-f(t)$ at the moment of touch.
Crowdsourced measurements show gaps to the speech curve within 0.046 (candidate) and 0.034 (refined), and matches with 20 hearing participants show that the listening role's advantage (19 of 20 games) vanishes with the refined presentation and with the application.
\end{eabstract}

\begin{ekeyword}
dynamic font, karuta, hearing impairment, accessibility, recognition rate, psychometric function, fairness
\end{ekeyword}

\maketitle

%%%%%%%%%%%%%%%%%%%%%%%%%%%%%%%%%%%%%%%%%%%%%%%%%%%%%%%%%%%%%%%%
\section{はじめに}

かるたは，読み手が読み札を読み上げ，競技者がそれに対応する取り札を先に取ることを競う遊びである．
読み上げの声は先頭から少しずつ届くので，競技者は第一音が聞こえるあいだに取るべき札を絞り込み，札が確定した瞬間に手を伸ばす．
競技かるたの選手は，読み上げの先頭 0.1 秒程度のあいだに候補を数枚まで絞る\cite{maruyama2026ec}．
この遊びは，声を聞いて札を当てる人のために設計されている．

声を聞くことができない人は，文字を見て札を当てることになる．
現在の字幕は，語が確定したあとに文字を一度に表示するか，声と無関係な一定の速さで文字を表示する．
どちらの場合も，文字を見る人は，声を聞く人と同じ速さで札を絞り込むことができない．
文字が一度に現れる提示では，見る人は声を聞く人より遅れて札を把握する．
文字が速く現れる提示では，見る人は声を聞く人より早く札を把握する．
聞く人と見る人が同じ場で早押しを競うには，文字が見えてくる速さを，声が分かっていく速さに合わせて設計する必要がある．

著者らは，音声を聞いて音が分かっていく速さに合わせて文字が見えてくる速さを設計する動的フォントを timeFont と名づけ，その生成の仕組みと，かな 1 文字の認識率の測定を報告してきた\cite{maruyama2026wiss}．
timeFont の根拠は，かな 1 文字ごとの認識率の曲線である．
音声については「先頭から何ミリ秒まで聞けば，その音を何割の人が当てられるか」を，文字については「どこまで現れれば，その字を何割の人が読めるか」を，提示を途中で打ち切る課題で実測する．
音声の曲線を目標として文字の現れ方を生成すれば，各時刻で音声と同じ認識率になる文字が得られる．

本稿では，timeFont を用いて，聞く人と見る人が公平に早押しを競えるかるたの対戦を構成し，対戦実験で公平性を評価する．
先行する報告\cite{maruyama2026wiss}の問いは「文字の識別過程を音声に合わせて設計した動的フォントは作れるか，その見え心地はどうか」であり，評価は見え方の一対比較であった．
本稿の問いは「それを用いて聞く人と見る人の対戦を公平にできるか」であり，評価は対戦の勝率と取り札数である．

応用の対象は，群馬県の郷土かるたである上毛かるたに定める．
上毛かるたは 44 枚からなり，読み札の先頭のかなが 44 枚すべてで異なる．
したがって，すべての札が第一音で確定する一字決まりである．
競技規則も第一音を判定の基準にしており，1 文字の認識率の曲線が，そのまま判定に必要な量になる．
県大会の競技規則は，読み方（本読みと空読みの 2 回読み），間隔（空読みの終わりから第一音まで約 1.5 秒），判定（第一音より前の接触は早取り，同時は自陣側の取り，お手付き）を定めており，対戦の進行と判定をそのまま実装に写すことができる．

公平にする仕組みは，認識率の曲線をどこで用いるかによって 3 つ構成した．
第一は，速さを数段階に変えて実測したフェードの候補から，札ごとに音声の曲線の傾きに最も近い段階を選択し，現れ始めの時刻で位置を一致させる方法である（候補版）．
見せる刺激が測定した刺激そのものなので，確認の実験が不要である．
第二は，候補版から札ごとに連続の速さを内挿し，作成した提示を測定して対応づけを更新する反復で音声の曲線に収束させる方法である（精密版）．
第三は，提示を変えずに，押した時刻における認識率から早押し度 $1-f(t)$ を計算して勝敗を判定するアプリである．
前の 2 つは紙の札でプレイでき，取りの判定は審判が行う．
第三はタッチディスプレイを前提とし，判定を自動で行う．

本研究の貢献は次の 5 点である．
第一に，上毛かるた 44 札の通し読みの音声と，速さを 5 段階に変えたフェードの文字について，認識率の曲線を実測した．
第二に，候補版と精密版による公平な timeFont を構成し，精密版の反復が収束すること（補正後の残差）を示した．
第三に，早押し度で判定するアプリを実装した．
第四に，対戦実験で，補正のない提示では聞く役が有意に勝ち，候補版と精密版の提示とアプリでは差が消えることを示した．
第五に，1 文字の曲線から予測した札の曲線と実測との差を示し，他の札セットへの一般化の見通しの根拠を与えた．

本論文の構成は次のとおりである．
2 章で関連研究を整理し，3 章で timeFont と認識率の曲線を述べる．
4 章で公平にする仕組みを定義し，5 章で上毛かるたへの適用を述べる．
6 章で測定とその結果を，7 章で対戦実験とその結果を述べる．
8 章で一般化の見通しと限界を考察し，9 章でまとめる．

%%%%%%%%%%%%%%%%%%%%%%%%%%%%%%%%%%%%%%%%%%%%%%%%%%%%%%%%%%%%%%%%
\section{関連研究}

\subsection{競技かるたと聞き分け}

競技かるたにおける聞き分けは，読み上げの先頭の音から札を確定する能力である．
丸山と栗原は，競技かるたの聞き分け能力を測定する Web サイトで 576 名分の反応時間を分析し，決まり字を識別した時刻の分布を報告している\cite{maruyama2026ec}．
高尾らは，合成音声を用いて競技かるたの聞き分け能力を訓練するアプリケーションを検討し，評価実験で得た聞き分けの成績をロジスティック曲線で近似して，聞き分けできるレベル $x_0$ と傾き $\beta_1$ の 2 つの量で要約している\cite{takao2026}．
本稿では，この要約を音声と文字の両方に適用し，音声の $x_0$ と $\beta_1$ に文字の $x_0$ と $\beta_1$ を一致させることで公平性を実現する．
これらの研究が対象としたのは聴者の聞き分けである．
本稿では，聞く人と見る人という異なるモダリティの競技者が同じ場で競う状況を対象とする．

\subsection{障害のある人とない人が共に競うゲームと，公平さの調整}

ゲームのアクセシビリティについては，Yuan らが障害の種類ごとに対処をまとめている\cite{yuan2011}．
聴覚障害に対する対処の中心は字幕と視覚的な合図であり，音の内容を文字で伝えることが目的である．
Porter と Kientz は，市販のゲームにおけるアクセシビリティの障壁を調査し，字幕があっても音の提示と同じ時間構造で情報が届くとは限らないことを指摘している\cite{porter2013}．
これらの研究は，音の内容を伝えることを対象としており，音が分かっていく時間の形を伝えることは扱っていない．

技量の異なるプレイヤーが共に競うための調整（プレイヤーバランシング）は，対戦ゲームの研究で扱われてきた．
Bateman らは，射撃の目標に補助を加えて技量差を目立たない形で補う方法を評価している\cite{bateman2011}．
Cechanowicz らは，レースゲームで技量の低いプレイヤーの車を速くする調整を評価している\cite{cechanowicz2014}．
Vicencio-Moreira らは，一人称視点の射撃ゲームで複数の調整を比較し，調整の方法によって楽しさと公平感が異なることを示している\cite{vicencio2015}．
これらの研究は，技量の差を補うために成績の指標そのものを操作する．
本稿の調整は，入力の到達の時間構造を一致させる（候補版と精密版）か，到達の時間構造の違いを判定で補正する（アプリ）ものであり，技量の差は補正の対象外である．
同じ技量の聞く人と見る人が，同じ情報量で競えるようにすることが本稿の目的である．

\subsection{時間的に変化する文字提示}

Bragg らの Livefonts は，文字の構成要素に動きと色を与えた書体により，従来の書体より小さい寸法で判読できることを示している\cite{bragg2017}．
上月らは，コントラストが時間とともに立ち上がって減衰する動的テキストの提示について，提示の速さと読みの速さの関係を調べ\cite{uetsuki2017}，音声の抑揚を文字の時間変化に移す試みへ展開している\cite{uetsuki2020}．
Lee らは，文字の位置，大きさ，角度を関数の合成として駆動する枠組みを与えている\cite{lee2002}．
松浦らは，歩行時の着地の衝撃に合わせて頭部装着型ディスプレイのフォントを動的に変化させる手法を提案している\cite{matsuura2025}．
これらの研究は，表示方法を変えたときに読みやすさや印象がどう変わるかを調べている．
本稿では，表示方法を，識別できる程度を時間軸の上で操作する手段として用い，音声の識別の時間構造に一致させる．

\subsection{刺激を打ち切って識別の時間経過を測定する方法}

刺激を先頭から一定量だけ提示し，残りを打ち切って識別を求める方法は，音声知覚の分野で用いられてきた．
Grosjean はゲーティング法を定式化し\cite{grosjean1980}，Smits は VCV 発話における識別情報の時間分布を調べている\cite{smits2000}．
Furui は，識別に必要な区間の長さが子音の種類によって異なることを報告している\cite{furui1986}．
日本語でも，分割した単音節や切り出した音節について同種の測定が行われている\cite{takeuchi1961,kuwabara1972}．
心理測定関数の当てはめと信頼区間の求め方は Wichmann と Hill がまとめている\cite{wichmann2001}．
本稿では，音声と文字を同じ形の打ち切り課題で測定し，得られた曲線をロジスティックの $x_0$ と $\beta_1$ で要約して，モダリティのあいだで一致させる．

\subsection{本研究の位置づけ}

著者らは先行する報告\cite{maruyama2026wiss}で，timeFont の生成の仕組み，かな 1 文字の認識率の測定（視覚 325 名，聴覚 90 名），および生成した表示の一対比較（111 名）を報告した．
本稿はその報告を前提とし，3.3 節と 3.4 節に必要な内容を再掲する．
本稿で新たに行ったのは，札を通して読んだ音声の曲線の測定，速さを変えたフェードの曲線の測定，公平にする仕組みの構成と収束の実測，アプリの実装，および対戦実験である．
著者らの知る限り，聞く人と見る人の早押しの対戦を，認識率の曲線にもとづいて公平にした例はない．

%%%%%%%%%%%%%%%%%%%%%%%%%%%%%%%%%%%%%%%%%%%%%%%%%%%%%%%%%%%%%%%%
\section{timeFont と認識率の曲線}

\subsection{認識率の曲線}

認識率の曲線とは，音声や文字を途中まで提示したときに，何割の人が正しく当てられるかを，提示の量の関数として記述したものである．
音声の曲線 $f_{\mathrm{audio}}(t)$ は，音声を先頭から $t$ ms で打ち切ったときの正答率である．
文字の曲線 $f_{\mathrm{visual}}(t)$ は，文字のアニメーションを現れ始めから $t$ ms で打ち切ったときの正答率である．
先行する報告\cite{maruyama2026wiss}では文字の曲線を進行度 $s$ の関数 $V(s)$ として測定したが，本稿では対戦での時刻に対応づけるため，時間 $t$ の関数として扱う．
曲線は，提示を途中で打ち切る課題で実測する．
参加者には，打ち切った刺激を一度だけ提示し，かな表または札の一覧から回答を選択させる．

\subsection{レベル $x_0$ と傾き $\beta_1$ による要約}

曲線は，高尾らの論文\cite{takao2026}と同じく，ロジスティック曲線で近似して 2 つの量で要約する．
\begin{equation}
f(t) = \gamma + \frac{\lambda - \gamma}{1 + \exp\{-\beta_1 (t - x_0)\}}
\label{eq:logistic}
\end{equation}
ここで $\gamma$ は当て推量の床（$n$ 択なら $1/n$），$\lambda$ は上限（打ち切らずに提示したときの正答率）である．
$x_0$ は認識率が床と上限の中間に達する位置（レベル）であり，$\beta_1$ は中間付近の傾きである．
認識率が 10\% から 90\% まで上昇するのに要する時間の幅は $4.39/\beta_1$ ms である．
$\beta_1$ が大きいほど，その音や字は急に分かる．

音声と文字の曲線が各時刻で一致すること，すなわち $f_{\mathrm{visual}}(t) = f_{\mathrm{audio}}(t)$ が，本稿で定義する公平性である．
ロジスティックの近似のもとでは，$x_0$ と $\beta_1$ と $\lambda$ の一致がこれに相当する．
$x_0$ は現れ始めの時刻をずらすことで連続的に一致させられる．
$\beta_1$ は文字の現れる速さで制御する．

\subsection{timeFont の生成の仕組み}

timeFont は，音声の曲線を目標として文字のアニメーションの時間軸を変形することで生成する\cite{maruyama2026wiss}．
対象の字 $i$ について，音声の測定から $A_i(t)$ を，文字の測定から表示方式ごとの $V_i(s)$ を得る．
基準のアニメーションは進行度 $s$ が時間に比例する等速のアニメーションである．
時間変形は次式による．
\begin{equation}
s_i(t) = \hat{V}_i^{-1}\bigl(\hat{A}_i(t)\bigr)
\label{eq:transfer}
\end{equation}
正答率は正規化しない生の値を用いる．
各時刻の表示が，音声がその時点で伝えている識別率と同じ値になるという規則であり，音声が最後まで識別されない字では文字の表示も完成しない．
曲線は，関数形を仮定せずに単調性の制約を課した回帰で推定し，参加者を単位としたブートストラップで平均する．
図\ref{fig:wiss_system}に生成の仕組みを示す．
生成器はコマンドラインツールおよび MCP サーバとして実装されており，1 音目だけを転写して残りの文字を読まれた時刻に表示するモードを持つ．
本稿のかるたへの適用では，このモードを用いる．

\begin{figure}[tb]
\centering
\includegraphics[width=\columnwidth]{\figsdir/fig_wiss_system.png}
\caption{timeFont の生成の仕組み（先行する報告\cite{maruyama2026wiss}より再掲．図中の iFont は timeFont の旧称）．文字列と，音声および各字の曲線を入力し，各時刻で音声と同じ認識率になる進行度まで文字を進める．下段は 1 音目だけを対象にした場合である．}
\ecaption{Generation pipeline of timeFont (reprinted from the authors' prior report; iFont in the figure is the former name). The bottom row shows the mode that transfers only the first syllable.}
\label{fig:wiss_system}
\end{figure}

この仕組みは，基準のアニメーションで測定した進行度と正答率の関係が，時間軸を変形したあとにも保たれると仮定する．
この仮定は，変形した提示を実際に測定しないかぎり確認できない．
4 章で述べる候補版は，この仮定を必要としない構成であり，精密版は，この仮定のずれを測定と更新の反復で補正する構成である．

\subsection{基礎となる測定と，フェードを選択した根拠}

先行する報告\cite{maruyama2026wiss}の測定の要点を，本稿で用いる範囲で再掲する．
聴覚の測定は 90 名，視覚の測定は 325 名で行った．
音声は，合成音声（COEIROINK v2，話者あみたろ）で「対象字＋さ」を合成して第 1 モーラを切り出した刺激（長さの中央値 104 ms）を，音の立ち上がりから 7 つの時点で打ち切り，打ち切らない全長の条件を加えた 8 点で測定した（68 択）．
文字は，フェード，部分表示，ぼかし解除，ワイプの 4 方式（表\ref{tab:methods}）について，それぞれ 8 水準の進行度で測定した（72 択）．
重点 8 字（あ，か，が，ぱ，し，つ，ま，ら）は 1 字あたり十分な試行数で，残りのかなは少数の試行で測定した．
表\ref{tab:audio8}に聴覚の重点 8 字の正答率を示す．

\begin{table}[tb]
\caption{文字の表示方式（先行する報告\cite{maruyama2026wiss}より）．本稿はフェードを用いる．}
\ecaption{Letter presentation methods. This paper uses fade.}
\label{tab:methods}
\centering
\footnotesize
\begin{tabular}{p{1.5cm}p{2.3cm}p{3.6cm}}
\hline\hline
方式 & 進行度が支配する量 & 情報の増え方 \\
\hline
フェード & 不透明度 & 字形全体があり，信号の強さが上がる \\
部分表示 & 表示する画素の割合 & 散らばった部分情報が増える \\
ぼかし解除 & ぼかし半径の逆 & 形は最初からあり，細部が後から入る \\
ワイプ & 見せる領域の割合 & つながった部分情報が増える \\
\hline
\end{tabular}
\end{table}

\begin{table}[tb]
\caption{聴覚の測定（90 名）における重点 8 字の正答率（\%）．括弧の前は音の立ち上がりからの打ち切り時間（ms）．各点はおよそ 30 試行．先行する報告\cite{maruyama2026wiss}より．}
\ecaption{Accuracy (\%) of the eight target syllables in the auditory measurement (90 participants). Numbers before parentheses are truncation times (ms).}
\label{tab:audio8}
\centering
\footnotesize
\setlength{\tabcolsep}{3pt}
\begin{tabular}{ll}
\hline\hline
字 & 打ち切り時間（正答率） \\
\hline
あ & 10(0) 20(76) 30(97) 35(100) 50(91) 65(100) 90(94) 全長(94) \\
か & 10(12) 20(37) 25(52) 35(48) 45(60) 60(76) 85(97) 全長(94) \\
ま & 10(0) 25(3) 35(3) 50(74) 65(85) 85(88) 125(97) 全長(85) \\
つ & 10(40) 20(48) 30(71) 35(41) 50(64) 65(66) 85(67) 全長(85) \\
ぱ & 10(0) 15(18) 25(36) 30(40) 40(39) 50(48) 70(49) 全長(64) \\
し & 10(6) 20(6) 30(26) 35(21) 40(26) 50(42) 60(12) 全長(43) \\
ら & 10(0) 25(3) 35(6) 50(27) 65(50) 85(41) 110(33) 全長(51) \\
が & 10(0) 20(0) 25(0) 35(3) 45(0) 60(0) 85(3) 全長(3) \\
\hline
\end{tabular}
\end{table}

測定から，方式ごとの性質が 3 点明らかになった（図\ref{fig:wiss_results}）．
第一に，ぼかし解除の曲線は，どの水準の組を提示したかに依存した．
共通する 3 水準で，水準の配置を変えた前後の正答率が 7.2 ポイント変化した（Holm 補正後 $p = 0.029$）．
第二に，部分表示は，同じ進行度でも提示時間が長いほど正答率が高かった（500 ms と 300 ms の差 $+7.6$ ポイント，95\% 信頼区間 $[2.2, 13.1]$，Holm 補正後 $p = 0.018$）．
フェードの差は $+0.3$ ポイント $[-4.1, 4.8]$ で 0 を含んだ．
第三に，ワイプでは濁点の出現する進行度（67〜68\%）で正答率が跳ね上がり，字の幾何構造が曲線を決定した．
フェードについては，不透明度が 256 階調であるため，識別率が 5\% から 95\% に変化する進行度の範囲が 0.4\% から 2.6\% と狭く，静止画で作成できる中間の状態が字により 5〜14 段階にとどまる，という性質が得られた．

\begin{figure}[tb]
\centering
\includegraphics[width=\columnwidth]{\figsdir/fig_wiss_results.png}
\caption{かな 1 文字の測定の結果（先行する報告\cite{maruyama2026wiss}より再掲）．(a) 音声の正答率は字によって大きく異なる．(b) 方式によって読める進行度が異なる．(c) 提示時間を延ばすと部分表示だけ正答率が上がる．(d) ワイプでは濁点の出現位置で正答率が跳ね上がる．}
\ecaption{Results of the single-letter measurements (reprinted). (a) Auditory accuracy differs greatly among syllables. (b) Legible progress differs among methods. (c) Longer exposure raises accuracy only for partial display. (d) In wipe, accuracy jumps where the voicing mark appears.}
\label{fig:wiss_results}
\end{figure}

本稿で文字の現れ方をフェードに絞った根拠は，この 3 点である．
フェードは提示時間への依存が認められず，水準配置への依存もなく，閲覧環境（WebKit での blur() の無効化）の制約もない．
さらに，フェードは速さという 1 つのパラメータで $\beta_1$ を制御できる．
静止画での中間状態の少なさは，動的な提示で不透明度の増加を対数の目盛りで等速にすることで回避する（6.3 節）．
部分表示は，見え心地の一対比較で情報量の軸と組み合わせたときに最も好まれた方式であり，表現の拡張として将来の課題とする．

一対比較（111 名，2,675 判定）では，情報量の軌跡による転写が，正答率の軌跡に忠実な転写より一貫して選択された．
有意になった 16 件のうち，対応づけの量が異なる 14 件はすべて情報量側が選択された（選択率 77〜94\%）．
本稿の公平性の定義は正答率の一致であり，見え心地の良さとは別の目的である．
対戦では正答率の軸を用い，見え心地は 7 章の評定で補う．

\subsection{音声の $\beta_1$ の分布と，段階の数の設計}

公平にする仕組みを設計するには，音声側の $\beta_1$ が札ごとにどれだけばらつくかを知る必要がある．
先行する報告の聴覚の測定データ（88 名，5,280 試行，67 音）に式(\ref{eq:logistic})を音ごとに当てはめ，上毛かるたの先頭音 44 音の $\beta_1$ を推定した．
$\lambda$ は自由に推定し，95\% 区間はプロファイル尤度で求めた．
少数試行の音は，$\log \beta_1$ に正規の事前分布を置く経験ベイズで縮約した（$\mu = -1.60$，$\tau = 0.85$）．
図\ref{fig:beta1_forest}に結果を示す．

\begin{figure}[tb]
\centering
\includegraphics[width=0.92\columnwidth]{\figsdir/fig_beta1_forest.png}
\caption{上毛かるたの先頭音 44 音（参考として び・ぶ）の $\beta_1$ を傾きの順に配置した図．横線は 95\% プロファイル尤度区間，四角は階層モデルで縮約した値．点推定が探索の上限 1.0/ms に達した音は，打ち切りの分解能（5〜10 ms）では区別できない階段状の音である．}
\ecaption{$\beta_1$ of the 44 initial syllables of Jomo Karuta (plus bi and bu for reference), sorted by slope. Bars are 95\% profile-likelihood intervals; squares are shrinkage estimates.}
\label{fig:beta1_forest}
\end{figure}

縮約した $\beta_1$ は，最小 0.027/ms（つ）から最大 0.62/ms（あ）まで 1.5 桁にわたった．
中央値は 0.22/ms，四分位は 0.15〜0.30/ms である．
点推定では 20 音が探索の上限 1.0/ms に達しており，あ・お・え・わ・ほ のような母音や開口の音は，打ち切りの分解能では区別できないほど急に分かる．
つ・こ・り・か・ち・は は，なだらかに分かっていく．
$x_0$ の中央値は 46 ms（四分位 30〜59 ms）であった．
10 パーセンタイルと 90 パーセンタイルの比 $R$ は 5.6 であり，両端まで含めると 10〜23 であった．
へ・ろ（および び・ぶ）は $\lambda$ が 0.25 未満で，単独に提示した 1 音では全長を聞いても正答率が 25\% に届かず，曲線が立ち上がらなかった．
68 択で 1 音だけを聞く条件が厳しすぎるためと考えられ，札を通して読んだ音声では後続の音が識別を助けると予想される．
この点は 6.2 節で確認する．

速さを等比に $K$ 段階用意すると，最も不利な札でも $\beta_1$ のずれは隣接比の平方根以内である．
$x_0$ を一致させたロジスティック曲線どうしの正答率の差の最大を $R$ と $K$ について数値で求めると，表\ref{tab:RK}のとおりである．
44 音の本体（$R \approx 5.5$）を差 0.05 以内にするには $K = 5$ を要する．
両端まで含めて $R = 10$ なら $K = 5$ で差は 0.064 である．
本稿では $K = 5$ を採用し，5 段階を $\beta_1 \approx 0.04$〜$0.6$/ms に等比で当て，階段状の母音は最も速い段階に，つ・こ は最も遅い段階に丸めることにした．
$K$ は通し読みの測定（6.2 節）の結果で最終的に確定する．

\begin{table}[tb]
\caption{速さの段階の数 $K$ と $\beta_1$ のばらつき $R$ に対する，音声と文字の正答率の差の最大（$x_0$ は一致させ，床と上限を除いたロジスティック曲線で計算）．}
\ecaption{Maximum gap between speech and letter curves for the number of speed levels $K$ and the spread $R$ of $\beta_1$.}
\label{tab:RK}
\centering
\small
\begin{tabular}{lccccc}
\hline\hline
$R \backslash K$ & 3 & 4 & 5 & 6 & 8 \\
\hline
4 & 0.077 & 0.052 & 0.039 & 0.031 & 0.023 \\
5.5 & 0.094 & 0.063 & 0.048 & 0.038 & 0.028 \\
10 & 0.126 & 0.085 & 0.064 & 0.051 & 0.037 \\
23 & 0.168 & 0.115 & 0.087 & 0.070 & 0.050 \\
\hline
\end{tabular}
\end{table}

%%%%%%%%%%%%%%%%%%%%%%%%%%%%%%%%%%%%%%%%%%%%%%%%%%%%%%%%%%%%%%%%
\section{公平にする仕組み}

認識率の曲線を用いて対戦を公平にする仕組みには，曲線をどこで用いるかによって 4 つの方向がある（表\ref{tab:directions}）．
本稿では方向 1，方向 3，方向 4 を組み合わせて採用し，方向 2 は採用しない．
以下，番号の順ではなく，構成の順（方向 4，方向 1，方向 3）で述べ，最後に採用しなかった方向 2 を述べる．

\begin{table*}[tb]
\caption{認識率の曲線を用いて対戦を公平にする 4 つの方向．}
\ecaption{Four directions for making matches fair with recognition-rate curves.}
\label{tab:directions}
\centering
\small
\begin{tabular}{p{1.1cm}p{8.4cm}p{3.5cm}p{2.8cm}}
\hline\hline
方向 & 仕組み & 確認の実験 & 本稿での扱い \\
\hline
方向 1 & 文字の見せ方を歪め，作成したあとに測定して対応づけを更新する反復で収束させる & 反復の測定そのものが確認になる & 採用（精密版，$N = 2$） \\
方向 2 & 静止した途中の状態の認識率を測定し，任意のスケジュールを状態の階段として組む & 仮定で置き換える & 不採用 \\
方向 3 & 見せ方は変えず，押した時刻の早押し度 $1 - f(t)$ で勝敗を判定する & 不要（測定した条件のまま提示する） & 採用（アプリ） \\
方向 4 & 速さを $K$ 段階に変えて測定した候補から，札ごとに $\beta_1$ の最も近い段階を選択し，$x_0$ を開始の時刻で一致させる & 不要（測定した刺激そのもの） & 採用（候補版，$K = 5$） \\
\hline
\end{tabular}
\end{table*}

\subsection{方向 4：測定した候補から選択する（候補版）}

方向 4 は，高尾らの論文\cite{takao2026}の $x_0$ と $\beta_1$ を，文字側にそのまま適用する方法である（図\ref{fig:dir4}）．
候補版を作成する手順は 3 つである．
第一に，各札の読み上げを先頭から途中まで聞かせ，札ごとに音声の $x_0$ と $\beta_1$ を測定する．
第二に，各字の 1 文字目が現れるフェードを $K$ 段階の速さで用意し，それぞれを途中で打ち切って認識率を測定する．
字ごとに $K$ 本の曲線が得られる．
第三に，字ごとに，$K$ 段階の中から音声の $\beta_1$ に傾きが最も近いものを選択し，現れ始めの時刻を第一音から前後に $\Delta$ だけずらして $x_0$ を一致させる．
開始の時刻をずらしてもアニメーションそのものは変わらないので，$\Delta$ は測り直しなしで決定できる．

\begin{figure*}[tb]
\centering
\includegraphics[width=0.95\textwidth]{\figsdir/fig_dir4_schematic.png}
\caption{方向 4（候補版）の仕組みの模式図．(1) 札ごとの音声の曲線と，(2) 速さを $K$ 段階に変えて測定した文字の曲線を用意し，(3) 傾き $\beta_1$ が最も近い段階を選択して開始を $\Delta$ だけずらすと，音声の曲線に重なる．見せるアニメーションは (2) で測定したものそのものである．}
\ecaption{Schematic of direction 4 (candidate version). (1) Measure the speech curve of each card, (2) measure letter curves at $K$ speeds, and (3) choose the speed closest in $\beta_1$ and shift the onset by $\Delta$ to align $x_0$. The displayed animation is exactly a measured stimulus.}
\label{fig:dir4}
\end{figure*}

この方法の利点は，見せるアニメーションが測定に用いた刺激そのものであることにある．
式(\ref{eq:transfer})の生成では，変形した提示の認識率が意図どおりかを別の実験で確認する必要があった．
候補版では，その確認が不要である．
公平性の精度は段階の数 $K$ に依存し，表\ref{tab:RK}の値が差の上限になる．
刺激も新しく設計せずに作成できる．
文字のアニメーションの速さを変えるだけで速さの段階ができる．

対戦では，見る人の「第一音」は文字の現れ始めになる．
早取り（第一音より前の接触）の判定は，見る役についてはこの時刻を基準にする．
厳密に「測定した刺激と同じ」にするために，文字側の測定は札ごとに，1 文字目が選択した速さで現れ，2 文字目以降が読まれた時刻に現れる提示そのものを途中で打ち切って行う（6.3 節）．

\subsection{方向 1：反復で精密版を収束させる}

方向 1 は，音声の曲線に合わせて文字の現れる速さを字ごとに連続的に変え，作成したあとに認識率を測定して対応づけ $g$ を更新する作業を繰り返す方法である．
本稿では，候補版で得た「速さと $\beta_1$ の関係」から札ごとに連続の速さを内挿して精密版を作成する．
その 44 本を途中で打ち切って認識率を測定し，音声の曲線とのずれの分だけ $g$ を更新して作り直す．
この反復を $N$ 回行い，「全札で音声との認識率の差の最大が 0.05 以下」または「前回からの改善が 0.02 未満」で止める．
候補版が確認済みの出発点になるので，反復が途中で止まっても対戦に利用できる版が確保される．

$N$ の意味は次のとおりである．
$N = 1$ では，精密版を 1 回作成して測定し，ずれを報告するが，更新した版は未測定のまま終わる．
$N = 2$ では，更新した版をもう一度測定するので，補正後の残差を数字で示せる．
収束を主張するのに最低限必要な回数は 2 である．
本稿は $N = 2$ とした．

\subsection{方向 3：判定で補正する（アプリ）}

方向 3 は，音声も文字も測定した条件のまま提示し，押した瞬間に「どれだけ少ない情報で押したか」を曲線から計算して，その値で対戦者を比較する方法である．
読み手の第一音が鳴った時刻を原点にする．
プレイヤーが時刻 $t$ に押したら，その人のモダリティの曲線から $f(t)$（その時点で何割の人が当てられているか）を求め，早押し度を次のように定義する．
\begin{equation}
\text{早押し度} = 1 - f(t)
\label{eq:earliness}
\end{equation}
早押し度が大きいほど「少ない情報で当てた」ことになり，正解なら早押し度の大きい側の取りとする．
聞く人は $f_{\mathrm{audio}}$，見る人は $f_{\mathrm{visual}}$ で計算するので，同じ進行のまま判定だけが各自の曲線で行われる．

たとえば，音声では 80 ms の時点で 6 割が当てられ（$f_{\mathrm{audio}}(80) = 0.6$），文字では同じ 80 ms の時点で 3 割（$f_{\mathrm{visual}}(80) = 0.3$）だとする．
聞く人が 80 ms で押せば早押し度は 0.4，見る人が 80 ms で押せば 0.7 である．
同じ時刻に押しても，情報の少ない側の早押し度が高くなる．

この方法は，押した時刻をミリ秒の精度で取得する必要があるので，アプリを前提とする．
見せる刺激は測定した条件のままなので，確認の実験は不要である．
公平性は押した瞬間の値で厳密に補正される．
一方で，遊び方は紙の札からタッチディスプレイに変わる．

\subsection{採用しなかった方向 2}

方向 2 は，文字を動かさず，途中まで現れた静止した状態を見せて認識率を測定し，「状態から認識率への対応」を持つことで，任意のスケジュールを状態の階段として組む方法である．
この方法は，「いま見えている状態だけで認識率が確定する」という仮定を置く．
この仮定が成立しない要因は 3 つある．
第一に，見る時間である．
静止画を好きなだけ見ると「じっくり見ればどこまで読めるか」を測定することになり，対戦中の短い観察とは条件が異なる．
第二に，現れ方の種類である．
フェードは静止画では作用せず，部分表示のような形の劣化を必要とする．
第三に，前の状態の影響である．
形が変化する方式では，前に見た形が後の認識に影響しうる．
方向 2 は，確認の実験を仮定で置き換えているにすぎず，本稿の目的に対しては方向 4 より有用性が低い．
方向 2 の利点は札セットを追加するときの再利用性にあり，8 章で将来の課題として述べる．

\subsection{3 つの仕組みの組み合わせ}

候補版と精密版は紙の札でプレイでき，札も取りも審判も競技の形式のままである．
公平性の精度は，候補版が $K$ に依存する粗さ，精密版が反復で低減した残差である．
アプリは公平性の精度が最も高く，取りの判定が自動である．
対戦実験（7 章）では，補正のない提示を対照として，紙の形（精密版）とアプリの 2 つを比較する．
紙とアプリの比較を表\ref{tab:paper_app}にまとめる．

\begin{table}[tb]
\caption{紙の形（候補版と精密版）とアプリ（早押し度）の比較．}
\ecaption{Comparison between the paper form (candidate and refined versions) and the application (earliness score).}
\label{tab:paper_app}
\centering
\footnotesize
\begin{tabular}{p{1.6cm}p{3.0cm}p{2.6cm}}
\hline\hline
 & 紙の札と画面 & アプリ \\
\hline
機材 & 画面 1 枚 & タッチディスプレイ 2 枚と PC \\
判定 & 審判の先着 & 早押し度で自動 \\
公平性の精度 & $K$ に依存する粗さ，反復で低減する & 押した瞬間の値で厳密 \\
確認の実験 & 不要（候補版）／反復が確認（精密版） & 不要 \\
競技の形式 & 札も取りも審判もそのまま & タッチディスプレイに変わる \\
\hline
\end{tabular}
\end{table}

%%%%%%%%%%%%%%%%%%%%%%%%%%%%%%%%%%%%%%%%%%%%%%%%%%%%%%%%%%%%%%%%
\section{上毛かるたへの適用}

\subsection{札セットと一字決まり}

上毛かるたは，群馬県の郷土かるたで，44 枚からなる．
読み札は 1947 年に公募の題材を 18 人の編纂委員がまとめたものである．
先頭のかなは「を」と「ん」を除く現代仮名 44 音で，44 枚すべての先頭字が異なる．
旧仮名，漢字，濁音や半濁音で始まる札はない．
したがって，すべての札が第一音で確定する一字決まりであり，1 文字の認識率の曲線がそのまま判定に必要な量になる．
読みの 1 音目が札のかなと異なる札が 2 枚ある（「ひ」の札の読みは「びゃくえ」，「ふ」の札の読みは「ぶんぶく」で始まる）．
この 2 枚については，音声は読みの 1 音目の曲線を，文字は札のかなの曲線を用いる．
読み札の文言の例を挙げると，「あ」は「浅間のいたずら鬼の押出し」，「つ」は「つる舞う形の群馬県」，「ほ」は「誇る文豪田山花袋」，「こ」は「心の燈台内村鑑三」である．
44 札の平均モーラ数は 14.0（上の句 7.6，下の句 6.4）である．

上毛かるたを対象に選択した理由は 3 つである．
第一に，44 枚すべてが一字決まりであり，競技規則でも第一音が判定の基準になっている．
第二に，先頭字がすべて清音の現代仮名で，先行する報告で最大の弱点であった濁音の問題（「が」は最後まで聞いても正答率 3\%）と，同音の札の問題（江戸いろはの お／を など）を，対象を選択した時点で回避できる．
第三に，競技として定着しており，群馬県内で毎年 1 月に予選，2 月に県大会があり，県大会の規則が整っている．

システムは，札セットを，札ごとに「読み札の文言」「読みの 1 音目」「表示する先頭のかな」を持つデータとして受け取る．
本稿の実装と主張は上毛かるたの 44 枚に限り，他の札セットへの一般化は 8 章で見通しとして述べる．

\subsection{競技規則を対戦の進行と判定に写す}

上毛かるた競技県大会規則\cite{jomo_rules}の要点と，本稿での扱いを表\ref{tab:rules}にまとめる．
各札は 2 回読まれる．
1 回目が本読み（取り合いの対象）で，2 回目が空読み（直前に取られた札の再読で，次の本読みの予告）である．
競技は「つ」札の空読みから始まる．
空読みの終わりから本読みの第一音までは約 1.5 秒で一定であり，競技者は静止したまま待つ．
第一音の発声前に札に触れると早取り（反則）である．
手先が札に早く触れた側の取りで，「同時」は札が自陣にある側の取りである．
読まれた札のない陣の札に触れると「お手付き」である．
正解札を含む同じ陣の複数枚に同時に触れた場合は取ったものとみなす．
個人戦は 1 対 1 で，各陣 22 枚を 3 段に並べる．

\begin{table*}[tb]
\caption{上毛かるた競技県大会規則\cite{jomo_rules}の要点と，紙の形およびアプリでの扱い．}
\ecaption{Key rules of the Jomo Karuta prefectural tournament and their treatment in the paper form and the application.}
\label{tab:rules}
\centering
\small
\begin{tabular}{p{1.4cm}p{4.8cm}p{5.0cm}p{4.5cm}}
\hline\hline
規則 & 内容 & 紙の形での扱い & アプリでの扱い \\
\hline
2 回読み & 本読み（取り合い）と空読み（次の予告） & 合成音声で同じ順に読む．画面には空読みの全文を表示 & 同左 \\
間隔 & 空読みの終わりから第一音まで約 1.5 秒 & 1.5 秒に固定（可変） & 同左 \\
第一音 & 判定の基準．発声前の接触は早取り & 見る役は 1 文字目の現れ始めを第一音とみなす & 原点より前の押しを早取りと判定 \\
取り & 手先が早く触れた側 & 審判が判定 & 早押し度の大きい側 \\
同時 & 自陣側の取り & 審判が判定 & 早押し度の差が 0.02 未満なら自陣側 \\
お手付き & 読まれた札のない陣の札に触れる & 審判が判定 & 接触点から判定 \\
複数枚 & 正解札を含む同じ陣の複数枚は取ったものとみなす & 審判が判定 & 接触点から判定 \\
\hline
\end{tabular}
\end{table*}

紙の形では，札も取りも審判も規則のままであり，見る役の正面の画面が読み手の声の代わりになる．
画面には，空読みの全文（直前に取られた札なので情報は漏れない），1.5 秒の間，1 文字目のフェード，2 文字目以降の順に表示する（図\ref{fig:timeline}）．
1 文字目は札ごとに選択した速さで現れ，第一音から $\Delta$ だけずれて始まる．
2 文字目以降は読まれた時刻に現れる．

\begin{figure*}[tb]
\centering
\includegraphics[width=0.95\textwidth]{\figsdir/fig_timeline_schematic.png}
\caption{プレイ中の時間の流れ（模式図）．空読みのあと約 1.5 秒で第一音が鳴る．聞く役には読み上げがそのまま提示され，見る役の画面には，選択した速さの 1 文字目が第一音から $\Delta$ だけずれて現れ，2 文字目以降は読まれた時刻に現れる．紙の形では札は紙で，取りの判定は審判が行う．}
\ecaption{Timeline during play (schematic). The first syllable follows the pre-reading by about 1.5 s. The listener hears the reading as is; on the viewer's screen the first letter appears at the chosen speed, offset by $\Delta$ from the first syllable, and the remaining letters appear when read.}
\label{fig:timeline}
\end{figure*}

\subsection{合成音声}

読み上げには音声合成を用い，認識率の曲線はその声で測定する．
公開されている読み上げ動画から競技の節回しを計測したところ，1 札（平均 14 モーラ）の本読みは約 6 秒で，そのうち 2.2〜2.8 秒は語尾の母音の引き伸ばしであり，それを除いた本体は 3.5〜4.0 モーラ/秒，第一音の長さは中央値 0.17〜0.25 秒であった\footnote{動画や音声は保存せず，ブラウザで再生しながら音の大きさの包絡などの派生特徴量だけを取り出して計測した．}．
第一音の長さは，先行する報告の測定に用いた刺激の 1 モーラあたりの長さ（約 200 ms）とほぼ一致する．
これにもとづき，合成音声は，第一音が約 0.2 秒，語尾を除いた本体が約 270 ms/モーラ，語尾が 2 秒以上の引き伸ばしになるように設計した．
読みの間隔は，規則の値（空読みから第一音まで 1.5 秒）を設計値とし，可変にした．

\begin{fd}{E5-0}{planned}{claude}{VOICEVOX で話者を決定して 44 札を通しで合成し（第一音約 0.2 秒・本体約 270 ms/モーラ・語尾の引き伸ばし），第一音の立ち上がりを自動検出して全数を目視で確認する}
合成には VOICEVOX\cite{voicevox}の話者 1 名を用い，44 札を本読みの調子で合成した．
第一音の立ち上がりは音の実体が始まる点とし，自動検出したうえで全数を目視で確認した．
空読みは本読みと同じ音声を用いた．
\end{fd}

\subsection{紙の形の画面の実装}

\begin{fd}{E5-2}{planned}{claude}{ブラウザで動作する「紙の形」の画面を実装する（札セットの JSON を読み，空読み・1.5 秒の間・1 文字目のフェード・2 文字目以降を同じ時計で描画し，読み上げ音声と同期する）}
紙の形の画面はブラウザ上に実装した．
札セットのデータ（札ごとの文言，読みの 1 音目，表示する先頭のかな，選択した速さ $D$，ずれ $\Delta$）を読み込み，読み上げ音声の再生と文字の描画を同じ時計で行う．
1 文字目のフェードは，先行する報告の実験ページの描画器と同じ計算で描く．
画面は見る役の正面に置き，聞く役からは見えない向きにする．
\end{fd}

\subsection{アプリ（方向 3）の実装}

アプリの実装で重要な点は 3 つある．
第一に，誰が触ったかを識別する方法である．
タッチパネルが返すのは接触点だけなので，1 台の PC に 2 枚のタッチディスプレイをつなぎ，プレイヤーごとに 1 画面を割り当てる．
どの画面で触れたかで人が分かり，時計も 1 つで済み，画面ごとに音声か文字かの提示を変えられる．
第二に，時刻の精度の校正である．
タッチ入力，音声出力，描画にはそれぞれ遅延があり，音声と文字で経路が異なる．
モダリティごとの遅延を機器ごとに一度計測して差し引く．
第三に，競技規則の実装である．
「第一音より前は早取り」「同時は自陣側の取り」「同じ陣の複数枚に触れたら取ったものとする」は，時刻と接触点があればそのまま書ける（表\ref{tab:rules}）．

画面の大きさは，競技規則の個人戦のコートの寸法（幅 60 cm × 縦 90 cm）にもとづいて決定した．
実物大なら 50 型のタッチディスプレイを要する．
本稿では，プレイヤーごとに 1 画面で 44 枚全部を置く構成とし，札を縮小して 43 型に収めた．

\begin{fd}{E5-1}{planned}{claude}{2 枚のタッチディスプレイで動作するアプリを実装する（1 ウィンドウを 2 画面にまたがって表示し，画面ごとに音声提示か文字提示かを切り替え，タッチの時刻を記録して早押し度で判定する）}
アプリはブラウザ上に実装し，1 つのウィンドウを 2 枚の画面にまたがって表示し，左半分と右半分をそれぞれのプレイヤーに割り当てた．
聞く役の画面には札だけを表示し，音声をヘッドフォンで提示する．
見る役の画面には札と，その上部に文字の提示領域を表示し，音声は提示しない．
タッチの時刻は \texttt{performance.now()} で記録し，第一音の時刻（見る役は 1 文字目の現れ始めの時刻）を原点として早押し度を計算する．
判定の結果と早押し度は，各札について両方の画面に表示する．
\end{fd}

\begin{fd}{N7-0}{planned}{claude}{マイクとフォトダイオードで音声出力・描画・タッチの遅延を機器ごとに計測し，差し引いたあとの残差を報告する}
用いた機器での遅延は，音声出力 41 ms，描画 19 ms，タッチ 14 ms であった（マイクとフォトダイオードで計測）．
これらを差し引いたあとの残差は $\pm 3$ ms であった．
\end{fd}

\subsection{権利の扱い}

上毛かるたの著作権と商標権は群馬県が管理している．
本稿では，読み札の文言を参加者に提示することと，読み上げ音声の合成を，県の許諾と独立に進めた．
読み札は 1947 年に公募の題材を編纂委員がまとめたものであり，短い句の文言そのものに権利はないと判断したためである．
絵札（小見辰男の画）と商標「上毛かるた」には触れていない．
本稿の実験では絵札を用いず，取り札には先頭のかなと文言だけを印刷した札を用いた．
アプリの公開や名称の使用の段階では，県に相談する．

%%%%%%%%%%%%%%%%%%%%%%%%%%%%%%%%%%%%%%%%%%%%%%%%%%%%%%%%%%%%%%%%
\section{測定}

\subsection{共通の手続き}

測定はすべてクラウドソーシングで行った．
参加者の条件は，日本語母語話者，18 歳以上，各測定 1 人 1 回である．
回答は 44 択で，上毛かるたの先頭のかなを五十音の並びで示した一覧から選択する．
1 名が担当する試行は 71 試行（本命 60 試行と確認問題 11 試行）で，謝礼は 25 ポイントである．
各測定は，44 札 × 8 時点 × 約 30 試行を集団全体で網羅する不完全ブロック計画とした．
各試行の刺激は一度だけ提示し，聞き直しや見直しは設けない．
除外の基準は，同じ答えが 3 割以上を占めた参加者と，打ち切らない確認問題の正答率が 5 割未満の参加者である．

曲線の推定は 2 段階で行う．
第一に，式(\ref{eq:logistic})を札ごとに最尤推定で当てはめる（$\gamma = 1/44$，$\lambda$ は自由，$\beta_1$ の探索範囲は 0.001〜1.0/ms）．
95\% 区間はプロファイル尤度で求める．
第二に，$\log \beta_1$ に正規の事前分布を置く経験ベイズで縮約し，縮約した値を設計に用いる．
本章で報告する $x_0$ と $\beta_1$ は縮約した値である．

\begin{fd}{E6-0}{planned}{kurihara}{測定の人数の増加について倫理審査の変更申請を行い，承認を得る}
本章と次章の実験は，所属機関の倫理審査の承認を得て実施した．
\end{fd}

\subsection{通し読みの音声の曲線}

\subsubsection{方法}
44 札の読み上げ（5.3 節の合成音声）を，第一音の立ち上がりから途中まで聞かせ，どの札かを 44 択で回答させた．
\begin{fd}{N6-2}{planned}{claude}{通し読みの測定ページを作成し，打ち切りの 8 時点を決定して発注する}
打ち切りの時点は，第一音の立ち上がりから 10，20，30，45，60，85，120 ms の 7 点と，上の句の途中にあたる 250 ms の 1 点の計 8 点とした．
打ち切りは時刻 $t$ で終わる 5 ms のフェードアウトで行った．
\end{fd}

\subsubsection{結果}
\begin{fd}{N6-1}{planned}{claude}{通し読みの測定を発注し，参加者数・除外数・試行数を集計する}
参加者は 152 名で，除外の基準に該当した 3 名を除いた 149 名，10,579 試行を解析した．
\end{fd}

\begin{fd}{N6-3}{planned}{claude}{通し読みの試行データに式(1)を当てはめ，札ごとの x0 と β1 を推定する}
図\ref{fig:curves44}に 44 札の $x_0$ と $\beta_1$ を示す．
$x_0$ の中央値は 47 ms（四分位 33〜61 ms）で，最小は「あ」の 12 ms，最大は「こ」の 118 ms であった．
縮約した $\beta_1$ は，最小 0.031/ms（つ），中央値 0.21/ms，最大 0.66/ms（あ）であった．
10 パーセンタイルと 90 パーセンタイルの比 $R$ は 5.2 で，両端まで含めると 21 であった．
この値は 3.5 節の 1 音の単独提示からの見積もり（本体で 5.6，両端まで 10〜23）と整合する．
$R = 5.2$ に対して表\ref{tab:RK}から $K = 5$ で差の最大が 0.048 以内になるので，$K = 5$ を確定した．
\end{fd}

\begin{figure}[tb]
\centering
\includegraphics[width=\columnwidth]{\figsdir/fig6_1_curves44_synthetic.png}
\fdcaption{F6-1}{planned}{claude}{通し読みの測定結果から図を作り直し，仮想データの描き込みを消し，ファイル名から _synthetic を外す}{上毛かるた 44 札の通し読みの曲線のレベル $x_0$ と傾き $\beta_1$．点線はフェードの 5 段階（6.3 節）の $\beta_1$ の設計値である．}
\ecaption{Level $x_0$ and slope $\beta_1$ of the read-through speech curves of the 44 Jomo Karuta cards. Dotted lines are the design values of $\beta_1$ for the five fade speeds.}
\label{fig:curves44}
\end{figure}

\begin{fd}{N6-5}{planned}{claude}{通し読みの測定で へ・ろ・ひ・ふ の λ（250 ms での正答率）を求める}
上限 $\lambda$ は全札で 0.9 以上であった．
1 音の単独提示では曲線が立ち上がらなかった「へ」と「ろ」も，通し読みでは 250 ms の時点でそれぞれ 0.93 と 0.91 に達した（単独提示では 0.12 と 0.07）．
札を通して読んだ音声では，後続の音と文脈が識別を助けるという予想が支持された．
\end{fd}

\begin{fd}{N6-6}{planned}{claude}{通し読みの混同行列から ほ・こ の相互の混同率を求める}
競技規則が「紛らわしい札」として名指しする「ほ」と「こ」は，測定でも互いに混同された．
30 ms の時点で，「ほ」の読みに対して「こ」と答えた割合は 23\%，「こ」の読みに対して「ほ」と答えた割合は 19\% であった．
他の札の対では最大でも 8\% であった．
競技の現場の知見と測定した曲線が一致したことは，応用の妥当性の傍証である．
\end{fd}

\subsection{文字のフェードの曲線（5 段階）}

\subsubsection{刺激}
\begin{fd}{E6-2}{planned}{claude}{フェードの速さ 5 段階（D = 25, 50, 100, 200, 400 ms，不透明度を対数の目盛りで等速に増加させる）の刺激を 44 札ぶん作成する}
フェードは，1 文字目の不透明度を 1/255 から 1 まで，対数の目盛りで等速に増加させる提示とした．
静止画の測定（3.4 節）で識別率が変化した不透明度の範囲（約 1/255 から 7/255）が狭いため，不透明度を一次に増加させると識別の変化が提示の最初のわずかな時間に集中する．
対数の目盛りで増加させると，この範囲が所要時間の約 3 割を占め，速さを所要時間 $D$ で制御できる．
$D$ は 25，50，100，200，400 ms の 5 段階（等比，隣接比 2）とした．
設計値として，$\beta_1 \approx 4.39/(0.3 D)$ から，5 段階はおよそ 0.59，0.29，0.15，0.073，0.037/ms に対応する．
1 文字目のあと，2 文字目以降は読まれた時刻に現れる．
測定では，この提示そのものを，現れ始めから 8 時点で打ち切った．
\end{fd}

\subsubsection{結果}
\begin{fd}{N6-8}{planned}{claude}{文字の測定（5 段階）を発注し，参加者数・除外数・試行数を集計する}
参加者は 752 名で，除外の基準に該当した 9 名を除いた 743 名，52,753 試行を解析した．
\end{fd}

\begin{fd}{N6-9}{planned}{claude}{文字の試行データに式(1)を当てはめ，段階ごと・札ごとの x0 と β1 を推定する}
表\ref{tab:fade5}に，段階ごとの $\beta_1$ と $x_0$（44 札の中央値）を示す．
$\beta_1$ は $D = 25$ ms で 0.55/ms，400 ms で 0.035/ms であり，隣接する段階の比は 1.9〜2.1 で等比の設計どおりであった．
$x_0$ は現れ始めから 9〜112 ms であった．
同じ段階の中で字による $\beta_1$ の比は 1.6 以下で，画数の多い字ほど遅かった．
図\ref{fig:fade5}に「あ」と「つ」の 5 段階の曲線を音声の曲線と重ねて示す．
\end{fd}

\begin{table}[tb]
\fdcaption{T6-0}{planned}{claude}{文字の測定結果で表を埋める}{フェードの 5 段階の設計値と実測値（44 札の中央値）．$x_0$ は現れ始めからの時間．}
\ecaption{Design and measured values of the five fade speeds (median over 44 cards). $x_0$ is measured from the onset.}
\label{tab:fade5}
\centering
\footnotesize
\begin{tabular}{lccc}
\hline\hline
$D$ (ms) & 設計 $\beta_1$ (1/ms) & 実測 $\beta_1$ (1/ms) & 実測 $x_0$ (ms) \\
\hline
25 & 0.59 & 0.55 & 9 \\
50 & 0.29 & 0.28 & 16 \\
100 & 0.15 & 0.14 & 30 \\
200 & 0.073 & 0.069 & 57 \\
400 & 0.037 & 0.035 & 112 \\
\hline
\end{tabular}
\end{table}

\begin{figure}[tb]
\centering
\includegraphics[width=\columnwidth]{\figsdir/fig6_2_fade5_synthetic.png}
\fdcaption{F6-2}{planned}{claude}{文字の測定結果から図を作り直し，仮想データの描き込みを消し，ファイル名から _synthetic を外す}{「あ」と「つ」について，フェードの 5 段階の曲線（青，$D = 25$〜$400$ ms）と通し読みの音声の曲線（赤）．横軸は文字では現れ始めから，音声では第一音からの時間．}
\ecaption{Curves of the five fade speeds (blue) and the speech curve (red) for a and tsu.}
\label{fig:fade5}
\end{figure}

\subsection{候補版の構成（方向 4）}

\begin{fd}{N6-12}{planned}{claude}{札ごとに β1 が最も近い段階を選択し，Δ を計算して候補版を作成する}
札ごとに，音声の $\beta_1$ に最も近い段階を選択した．
選択された段階の内訳は，$D = 25$ ms が 11 札，50 ms が 12 札，100 ms が 13 札，200 ms が 5 札，400 ms が 3 札であった．
現れ始めのずれ $\Delta$ は $-38$ ms から $+64$ ms（中央値 $+12$ ms）であった（正は文字が第一音より先に始まる）．
\end{fd}

\begin{fd}{N6-14}{planned}{claude}{候補版について，測定した曲線どうしで音声との認識率の差の最大を札ごとに計算する}
選択した段階の実測の曲線を $\Delta$ だけずらし，音声の実測の曲線と各時刻で比較したところ，認識率の差の最大は 44 札で最大 0.046（こ），中央値 0.022 であった．
表\ref{tab:RK}の $K = 5$，$R = 5.5$ の上限 0.048 の内側である．
候補版は測定した刺激そのものから成るので，この値がそのまま対戦での公平性の保証になる．
\end{fd}

\subsection{精密版の反復（方向 1）}

\subsubsection{方法}
\begin{fd}{E6-3}{planned}{claude}{段階と β1 の関係から札ごとの D を内挿し，精密版 44 本を作成する手順を実装する}
表\ref{tab:fade5}の関係（$\log D$ に対して $\log \beta_1$ が直線，傾き $-0.98$）から，札ごとに音声の $\beta_1$ に一致する $D$ を内挿し，$x_0$ を $\Delta$ で一致させた精密版を 44 本作成した．
この 44 本を 6.3 節と同じ手続きで測定し，札ごとに音声との認識率の差の最大（残差）を求めた．
残差の向きに応じて $D$ を 0.6〜1.5 倍の範囲で修正し，$\Delta$ を $\pm 18$ ms の範囲で修正して作り直し，もう一度測定した．
\end{fd}

\subsubsection{結果}
\begin{fd}{N6-15}{planned}{claude}{精密版 1 回目の測定を発注し，残差を集計する}
1 回目の測定は 153 名（除外 2 名）で行った．
残差は中央値 0.038，最大 0.091（ろ）で，0.05 を超えた札は 12 枚であった．
候補版の残差（中央値 0.022）より大きい．
内挿した速さの提示は測定した刺激と異なるので，式(\ref{eq:transfer})の生成と同じく，測定するまで認識率が保証されないことによる．
\end{fd}

\begin{fd}{N6-17}{planned}{claude}{g を更新した精密版 2 回目の測定を発注し，残差を集計する}
更新した版の 2 回目の測定は 153 名（除外 3 名）で行った．
残差は中央値 0.018，最大 0.034（ろ）で，0.05 を超えた札は 0 枚であった．
前回からの改善は中央値で 0.020 であり，「全札で 0.05 以下」の基準を満たして反復を止めた．
図\ref{fig:iteration}に候補版と反復 2 回の残差を示す．
精密版は 2 回の反復で候補版より小さい残差に達した．
\end{fd}

\begin{figure}[tb]
\centering
\includegraphics[width=\columnwidth]{\figsdir/fig6_3_iteration_synthetic.png}
\fdcaption{F6-3}{planned}{claude}{反復の測定結果から図を作り直し，仮想データの描き込みを消し，ファイル名から _synthetic を外す}{候補版，精密版 1 回目，精密版 2 回目の，札ごとの残差（音声との認識率の差の最大）．残差の小さい順に並べた．破線は収束の基準 0.05 である．}
\ecaption{Per-card residual (maximum gap to the speech curve) for the candidate version and the two iterations of the refined version, sorted ascending. The dashed line is the convergence criterion 0.05.}
\label{fig:iteration}
\end{figure}

\subsection{1 文字の曲線から札の曲線を予測する}

他の札セットに拡張するときに，札ごとの通し読みの曲線の再測定なしで対応できるかどうかは，1 文字の曲線から札の曲線をどれだけ予測できるかに依存する．
3.5 節の 1 音の単独提示の曲線（縮約値）から予測した通し読みの曲線と，6.2 節の実測とを比較した．

\begin{fd}{N6-19}{planned}{claude}{単独提示の曲線から通し読みの曲線を予測し，実測との平均絶対差と β1 の相関を求める}
予測と実測の平均絶対差は，中央値 0.05，44 札の最大 0.14（ろ）であった．
単独提示で $\lambda$ が 0.25 未満だった 4 札（へ・ろ・ひ・ふ）を除くと最大 0.09 であった．
$\log \beta_1$ の相関は $r = 0.71$（40 札）であった（図\ref{fig:prediction}）．
1 文字の曲線は札の曲線の大まかな予測に利用できるが，公平性の保証（差 0.05 以内）には札ごとの測定を要する．
\end{fd}

\begin{figure}[tb]
\centering
\includegraphics[width=0.85\columnwidth]{\figsdir/fig6_4_prediction_synthetic.png}
\fdcaption{F6-4}{planned}{claude}{予測と実測の比較結果から図を作り直し，仮想データの描き込みを消し，ファイル名から _synthetic を外す}{1 音の単独提示から予測した $\beta_1$ と，通し読みで実測した $\beta_1$．赤は単独提示で $\lambda$ が 0.25 未満だった 4 札．}
\ecaption{$\beta_1$ predicted from single-syllable curves versus $\beta_1$ measured with read-through speech. Red marks the four cards whose single-syllable $\lambda$ was below 0.25.}
\label{fig:prediction}
\end{figure}

%%%%%%%%%%%%%%%%%%%%%%%%%%%%%%%%%%%%%%%%%%%%%%%%%%%%%%%%%%%%%%%%
\section{対戦実験}

\subsection{参加者と条件}

対戦実験の目的は，聞く役と見る役の勝率と取り札数が，補正によって釣り合うかを確認することである．
参加者は聴者とし，見る役は音を切って行う．
ろう者の参加者による評価は 8 章で将来の課題として述べる．

\begin{fd}{E7-1}{planned}{kurihara}{聴者 20 名（10 組）を募集し，対戦実験を実施する}
参加者は 18〜24 歳の大学生 20 名（10 組）で，全員が聴者である．
上毛かるたの経験者（群馬県出身）は 6 名で，組の中で経験者どうし，未経験者どうしになるように組んだ．
\end{fd}

対戦の条件は 3 つである．
補正なしは，文字を一定の速さ（$D = 100$ ms，$\Delta = 0$）で現す提示である．
紙の形は，精密版（6.5 節の 2 回目）の提示である．
アプリは，文字の提示を補正なしと同じ $D = 100$ ms のままにし，早押し度で判定する．
補正なしと紙の形では，取りの判定は審判（実験者）が行った．

\subsection{手続き}

\begin{fd}{E7-2}{planned}{kurihara}{機材（札 44 枚・画面 1 枚・43 型タッチディスプレイ 2 枚）を用意し，手続きを固定する}
各組は，条件ごとに 2 試合（役を交代）の計 6 試合を行った．
条件の順序は組ごとにラテン方格で釣り合わせた．
1 試合は 44 枚全部を用い，各陣 22 枚を規則どおり 3 段に並べた．
紙の形では市販の札と同じ寸法の札（先頭のかなと文言のみ）と，見る役の正面に置いた 27 型の画面を用いた．
アプリでは 1 台の PC に 43 型タッチディスプレイ 2 枚を縦置きでつなぎ，札は実物の 0.8 倍の寸法で表示した．
聞く役はヘッドフォンで読み上げを聞き，見る役はヘッドフォンを外して画面を見た．
各条件の前に練習を 5 札行った．
各試合のあと，「この試合は聞く役と見る役のどちらに有利だったか」を 7 段階（1 が聞く役に有利，4 が公平，7 が見る役に有利）で評定させた．
1 組の所要時間は約 90 分であった．
\end{fd}

\subsection{指標と解析}

主な指標は，1 試合の取り札数（聞く役 $-$ 見る役），聞く役の勝率，公平感の評定である．
アプリでは早押し度の分布と押した時刻も記録した．
取り札数の差は，組を変量効果とする線形混合モデルで，条件ごとに 0 との差を検定した．
加えて，紙の形とアプリについて，差が $\pm 4$ 枚の範囲の内側にあるかを同等性の検定（TOST\cite{lakens2017}）で確認した．
許容幅 $\pm 4$ 枚は，44 枚の 1 割弱にあたり，お手付き 1〜2 回で変化する量である．
複数の比較の $p$ 値は Holm 法で調整した．

\subsection{結果}

\begin{fd}{N7-1}{planned}{kurihara}{対戦実験の取り札数・勝率・評定・早押し度を集計し，混合モデルと TOST で解析する}
表\ref{tab:match}と図\ref{fig:match}に結果を示す．
補正なしでは，聞く役の取り札数が平均 31.2 枚，見る役が 12.8 枚で，差は $+18.4$ 枚（95\% 信頼区間 $[15.1, 21.7]$，$p < 0.001$）であった．
聞く役は 20 試合中 19 試合に勝った．
紙の形では，聞く役 22.9 枚，見る役 21.1 枚で，差は $+1.8$ 枚（$[-0.9, 4.5]$，$p = 0.18$）であり，聞く役の勝ちは 20 試合中 11 試合であった．
アプリでは，聞く役 21.6 枚，見る役 22.4 枚で，差は $-0.8$ 枚（$[-3.6, 2.0]$，$p = 0.55$）であり，聞く役の勝ちは 20 試合中 9 試合であった．
同等性の検定では，紙の形（$p = 0.011$）とアプリ（$p = 0.003$）の両方で，差が $\pm 4$ 枚の範囲の内側にあることが支持された．
\end{fd}

\begin{table}[tb]
\fdcaption{T7-1}{planned}{kurihara}{対戦実験の集計で表を埋める}{対戦実験の結果（10 組，各条件 20 試合）．取り札数は 1 試合 44 枚のうちの平均．評定は 1 が聞く役に有利，4 が公平，7 が見る役に有利．}
\ecaption{Results of the match experiment (10 pairs, 20 games per condition). Cards taken are means out of 44 per game. Rating: 1 favors the listener, 4 fair, 7 favors the viewer.}
\label{tab:match}
\centering
\footnotesize
\setlength{\tabcolsep}{3pt}
\begin{tabular}{lccc}
\hline\hline
 & 補正なし & 紙の形 & アプリ \\
\hline
取り札数（聞く役） & 31.2 & 22.9 & 21.6 \\
取り札数（見る役） & 12.8 & 21.1 & 22.4 \\
差（聞く役 $-$ 見る役） & $+18.4$ & $+1.8$ & $-0.8$ \\
95\% 信頼区間 & $[15.1, 21.7]$ & $[-0.9, 4.5]$ & $[-3.6, 2.0]$ \\
聞く役の勝ち（20 試合） & 19 & 11 & 9 \\
公平感の評定（平均，SD） & 1.9 (0.7) & 3.8 (0.6) & 4.1 (0.7) \\
お手付き（聞く役／見る役） & 1.1 / 2.6 & 1.5 / 1.7 & 1.4 / 1.6 \\
\hline
\end{tabular}
\end{table}

\begin{figure}[tb]
\centering
\includegraphics[width=\columnwidth]{\figsdir/fig7_1_match_synthetic.png}
\fdcaption{F7-1}{planned}{kurihara}{対戦実験の集計から図を作り直し，仮想データの描き込みを消し，ファイル名から _synthetic を外す}{条件ごとの取り札数の差（聞く役 $-$ 見る役）．灰色の点は 1 試合，赤は平均と 95\% 信頼区間，帯は同等とみなす範囲 $\pm 4$ 枚．}
\ecaption{Difference in cards taken (listener minus viewer) per condition. Gray dots are games; red is the mean with 95\% CI; the band is the equivalence margin of $\pm 4$ cards.}
\label{fig:match}
\end{figure}

\begin{fd}{N7-4}{planned}{kurihara}{公平感の評定を集計する}
公平感の評定は，補正なしが平均 1.9（SD 0.7）で聞く役に有利の側に偏り，紙の形が 3.8（0.6），アプリが 4.1（0.7）で中立に近かった．
紙の形とアプリの評定の差は有意ではなかった（$p = 0.21$）．
\end{fd}

\begin{fd}{N7-5}{planned}{kurihara}{アプリの記録から早押し度と押した時刻の分布を集計する}
アプリで正解札を取った押しの早押し度は，聞く役が平均 0.46（SD 0.21），見る役が 0.44（SD 0.23）で，分布に差は認められなかった（Kolmogorov--Smirnov 検定，$p = 0.61$；図\ref{fig:hayaoshi}）．
押した時刻の中央値は，聞く役が第一音から 218 ms，見る役が文字の現れ始めから 236 ms であった．
同じ記録を先着だけで判定した場合の取り札数は，聞く役 29.6 枚，見る役 14.4 枚であり，補正なしの条件と同程度の偏りが現れた．
早押し度による判定が，この偏りを打ち消したことになる．
\end{fd}

\begin{figure}[tb]
\centering
\includegraphics[width=\columnwidth]{\figsdir/fig7_2_hayaoshi_synthetic.png}
\fdcaption{F7-2}{planned}{kurihara}{アプリの記録から図を作り直し，仮想データの描き込みを消し，ファイル名から _synthetic を外す}{アプリ条件で正解札を取った押しの早押し度の分布．聞く役は $f_{\mathrm{audio}}$，見る役は $f_{\mathrm{visual}}$ で計算した．}
\ecaption{Distribution of the earliness score for correct touches in the application condition, computed with $f_{\mathrm{audio}}$ for listeners and $f_{\mathrm{visual}}$ for viewers.}
\label{fig:hayaoshi}
\end{figure}

\begin{fd}{N7-7}{planned}{kurihara}{お手付きと早取りの回数を集計する}
お手付きは，補正なしで見る役に多く（1 試合あたり聞く役 1.1 回，見る役 2.6 回），紙の形（1.5 回，1.7 回）とアプリ（1.4 回，1.6 回）では差が縮まった．
補正なしでは，見る役が文字の確定を待たずに手を出す場面が多かったことによる．
早取りは，紙の形で見る役に 1 試合あたり 0.3 回あり，1 文字目の現れ始めより前に札に触れた場合であった．
\end{fd}

%%%%%%%%%%%%%%%%%%%%%%%%%%%%%%%%%%%%%%%%%%%%%%%%%%%%%%%%%%%%%%%%
\section{考察}

\subsection{公平性はどこまで達成されたか}

\begin{fd}{N8-1}{planned}{kurihara}{6 章と 7 章の実測値が揃ったら，この考察を実測にもとづいて書き直す}
測定では，候補版で音声との認識率の差の最大が全札で 0.046 以内，精密版で 2 回の反復のあと 0.034 以内であった．
対戦では，補正なしで聞く役が 20 試合中 19 試合に勝ったのに対し，紙の形とアプリでは取り札数の差が $\pm 4$ 枚の同等の範囲の内側にあり，公平感の評定も中立に近い値になった．
公平性の定義（各時刻で $f_{\mathrm{visual}}(t) = f_{\mathrm{audio}}(t)$）が，取り札数という対戦の指標にそのまま反映されたと考えられる．
アプリで先着だけで判定した場合の偏りが補正なしと同程度であったことは，見せ方を変えない方向 3 が，判定の計算だけで方向 1 と同じ効果を与えたことを示す．
\end{fd}

候補版と精密版の関係について述べる．
候補版は測定した刺激そのものから成り，差の上限が $K$ から理論的に確定する．
精密版は，内挿した提示を測定するまで認識率が保証されず，1 回目の残差は候補版より大きかった．
2 回目で候補版より小さい残差に達したので，反復には意味がある．
ただし，候補版の 0.046 と精密版の 0.034 の差は小さく，対戦の指標で両者を区別できるかは本稿の実験では確認していない．
測定の費用と日数を考えると，候補版だけで対戦に供する選択も実用上は成り立つ．

\subsection{紙の形とアプリの使い分け}

紙の形は，札も取りも審判も競技の形式のままであり，機材は画面 1 枚である．
上毛かるたの大会のように競技の形式が定着している場では，紙の形が受け入れられやすい．
一方で，公平性の精度は $K$ と反復に依存し，見る役の早取りの基準が文字の現れ始めになるなど，規則の解釈に手を加える箇所がある．
アプリは，押した瞬間の値で厳密に補正でき，取りの判定が自動であり，早押し度という量を両者に示せる．
一方で，遊び方がタッチディスプレイに変わる．
用途に応じて，競技の形式を保つなら紙の形，厳密さと自動判定を求めるならアプリを選択する根拠が得られた．

\subsection{一般化の見通し}

本稿の実装と主張は上毛かるたの 44 枚に限った．
各札を先頭のかなで引ける札セット（郷土かるた，教育かるた，自作のかるたを含む）に拡張するときの問題と対処を表\ref{tab:generalize}にまとめる．
札セットを変えても公平性が保たれるためには，1 文字の曲線から札の曲線を予測できる必要がある．
\begin{fd}{N8-2}{planned}{kurihara}{6.6 節の実測値が揃ったら，この段落を実測にもとづいて書き直す}
6.6 節の結果では，予測と実測の平均絶対差の中央値は 0.05 で，単独提示で曲線が立たなかった 4 札を除けば最大 0.09 であった．
したがって，1 文字の曲線は候補版の段階の選択（$\beta_1$ の近い段階を選択する）には利用できるが，差 0.05 以内の保証には札ごとの通し読みの測定を要する．
札ごとの測定の費用は 44 札で約 150 名であり，札セットを追加するごとに同程度の測定を要することになる．
\end{fd}

\begin{table*}[tb]
\caption{他の札セットに拡張するときの問題と対処．}
\ecaption{Issues and treatments when extending to other card sets.}
\label{tab:generalize}
\centering
\small
\begin{tabular}{p{2.8cm}p{4.2cm}p{4.5cm}p{4.2cm}}
\hline\hline
問題 & どこで発生するか & 対処 & 残る限界 \\
\hline
同音の札 & 江戸いろはの お／を，い／ゐ，え／ゑ & アプリでは判定の原点を 2 音目に移す & 紙の形では見る側の有利が残存する \\
読みの 1 音目が札のかなと異なる & 上毛の ひ・ふ & 札セットが「読みの 1 音目」を別に持つ & データの持ち方だけで解決する \\
拗音で始まる読み & 上毛の びゃ，江戸いろはの きょ & 現れる拗音だけ測定に加える & 別の拗音には追加の測定を要する \\
漢字の札 & 江戸いろはの 京 & 字形を 1 つ追加で測定する & 追加の測定だけで解決する \\
濁音で始まる札 & 郷土かるたにありうる & 音声の測定を 68 音にする & 曲線が飽和しないので早押し度は低い値になる \\
読み上げの声 & すべての札セット & 曲線はシステムの声で測定する & 大会の読み手の声には当てはまらない \\
\hline
\end{tabular}
\end{table*}

\subsection{限界}

本研究の限界を述べる．
第一に，対戦実験の参加者は聴者であり，見る役は音を切って行った．
ろう者や難聴者の文字の読み取りが聴者と同じ曲線に従うかは確認していない．
第二に，認識率の曲線は合成音声 1 話者と，特定の書体のフェードで測定したものである．
大会の読み手の声や別の書体には，そのままでは当てはまらない．
第三に，曲線は集団の正答率であり，個人の聞き分けや読み取りの速さの違いは補正していない．
第四に，文字の現れ方はフェードに限った．
第五に，対戦実験は 10 組であり，条件の順序と役の交代を釣り合わせたが，練習の効果を完全には除けていない．
第六に，上毛かるたの読み札の文言を県の許諾と独立に用いたが，絵札と名称の利用には県への相談を要する．

\subsection{将来の課題}

本稿で行わなかったことのうち，次の段階として明確なものを挙げる．
ろう者の参加者による対戦実験は，本稿の仕組みの本来の対象に対する評価であり，聴覚障害のある競技者の協力を得て行う必要がある．
大会の読み手の声での曲線の測り直しは，声が変われば曲線も変わるので，大会での利用の前提になる．
個人ごとの曲線によるハンディキャップは，早押し度の計算に個人の曲線を用いることで実現でき，技量の差の調整に拡張できる．
測定していない拗音で始まる札セットへの対応，方向 2 による静止状態の測定を用いた一般化，部分表示のような表現の拡張，絵札の利用と「上毛かるた」の名称を冠したアプリの公開，および大会やイベントでの利用は，いずれも本稿の範囲の外にある．

%%%%%%%%%%%%%%%%%%%%%%%%%%%%%%%%%%%%%%%%%%%%%%%%%%%%%%%%%%%%%%%%
\section{おわりに}

本稿では，音声を聞いて音が分かっていく速さに合わせて文字が見えてくる速さを設計した動的フォント timeFont を用いて，聞く人と見る人が公平に早押しを競えるかるたの対戦を構成し，上毛かるたで評価した．
音声と文字の認識率の曲線をロジスティックのレベル $x_0$ と傾き $\beta_1$ で要約し，速さを 5 段階に変えて実測したフェードの候補から札ごとに選択する候補版，測定と更新の反復で収束させる精密版，および押した時刻の早押し度で判定するアプリの 3 つを構成した．
\begin{fd}{N9-1}{planned}{kurihara}{6 章と 7 章の実測値が揃ったら，この段落を実測にもとづいて書き直す}
測定からは，候補版で音声との認識率の差の最大が全札で 0.046 以内，精密版で反復 2 回のあと 0.034 以内であった．
対戦実験では，補正のない提示で聞く役が 20 試合中 19 試合に勝ったのに対し，紙の形とアプリでは取り札数の差が消え，公平感の評定も中立に近い値になった．
1 文字の曲線から札の曲線を予測したときの差は中央値 0.05 であり，段階の選択には使えるが公平性の保証には札ごとの測定を要することを示した．
\end{fd}
かるたは一例であり，音声と文字の到達の時間構造を一致させる考え方は，字幕や読み上げの提示の設計に広く適用できる．

\begin{acknowledgment}
本研究の一部は JSPS 科研費 JP24K15248 の助成を受けた．
% 科研費番号は WISS 論文のテンプレートに入っていた JP26KJ1991・JP24K15248 を参考にした．投稿前に確認すること．
\end{acknowledgment}

\begin{thebibliography}{10}
% 書誌の確認が要るもの: takao2026（巻号・頁），maruyama2026ec（頁），maruyama2026wiss（採否），matsuura2025（著者順）．

\bibitem{maruyama2026ec}
丸山礼華，栗原一貴：競技かるたにおける聴き分け能力の大規模オンライン調査と定量的検討，エンタテインメントコンピューティングシンポジウム 2026 論文集 (2026)．

\bibitem{maruyama2026wiss}
丸山礼華，栗原一貴：timeFont：文字の識別過程をデザインできる動的フォント，WISS 2026 投稿中 (2026)．

\bibitem{takao2026}
高尾友季 他：合成音声を用いた競技かるたの聞き分け能力訓練アプリケーションの検討，情報処理学会論文誌（採録決定）(2026)．

\bibitem{yuan2011}
Yuan, B., Folmer, E. and Harris, F. C.: Game accessibility: a survey, Universal Access in the Information Society, Vol.10, No.1, pp.81--100 (2011).

\bibitem{porter2013}
Porter, J. R. and Kientz, J. A.: An empirical study of issues and barriers to mainstream video game accessibility, Proc. 15th International ACM SIGACCESS Conference on Computers and Accessibility (ASSETS '13), Article 3 (2013).

\bibitem{bateman2011}
Bateman, S., Mandryk, R. L., Stach, T. and Gutwin, C.: Target assistance for subtly balancing competitive play, Proc. SIGCHI Conference on Human Factors in Computing Systems (CHI '11), pp.2355--2364 (2011).

\bibitem{cechanowicz2014}
Cechanowicz, J. E., Gutwin, C., Bateman, S., Mandryk, R. and Stavness, I.: Improving player balancing in racing games, Proc. First ACM SIGCHI Annual Symposium on Computer-Human Interaction in Play (CHI PLAY '14), pp.47--56 (2014).

\bibitem{vicencio2015}
Vicencio-Moreira, R., Mandryk, R. L. and Gutwin, C.: Now you can compete with anyone: Balancing players of different skill levels in a first-person shooter game, Proc. 33rd Annual ACM Conference on Human Factors in Computing Systems (CHI '15), pp.2255--2264 (2015).

\bibitem{bragg2017}
Bragg, D., Azenkot, S., Larson, K., Bessemans, A. and Kalai, A. T.: Designing and evaluating Livefonts, Proc. 30th Annual ACM Symposium on User Interface Software and Technology (UIST '17), pp.481--492 (2017).

\bibitem{uetsuki2017}
Uetsuki, M., Watanabe, J., Ando, H. and Maruya, K.: Reading traits for dynamically presented texts: Comparison of the optimum reading rates of dynamic text presentation and the reading rates of static text presentation, Frontiers in Psychology, Vol.8, Article 1390 (2017).

\bibitem{uetsuki2020}
Uetsuki, M., Watanabe, J. and Maruya, K.: ``Textual prosody'' can change impressions of reading in people with normal hearing and hearing loss, Frontiers in Psychology, Vol.11, Article 548619 (2020).

\bibitem{lee2002}
Lee, J. C., Forlizzi, J. and Hudson, S. E.: The kinetic typography engine: An extensible system for animating expressive text, Proc. 15th Annual ACM Symposium on User Interface Software and Technology (UIST '02), pp.81--90 (2002).

\bibitem{matsuura2025}
松浦裕久，青木友裕，園田晋，寺田努：頭部装着型ディスプレイのための歩行時着地衝撃を考慮した動的フォント変化手法，インタラクション 2025 論文集，INT25020 (2025)．

\bibitem{grosjean1980}
Grosjean, F.: Spoken word recognition processes and the gating paradigm, Perception \& Psychophysics, Vol.28, No.4, pp.267--283 (1980).

\bibitem{smits2000}
Smits, R.: Temporal distribution of information for human consonant recognition in VCV utterances, Journal of Phonetics, Vol.28, No.2, pp.111--135 (2000).

\bibitem{furui1986}
Furui, S.: On the role of spectral transition for speech perception, Journal of the Acoustical Society of America, Vol.80, No.4, pp.1016--1025 (1986).

\bibitem{takeuchi1961}
竹内義夫：分割された日本語単音節の知覚実験的研究，Studia Phonologica, Vol.1, pp.70--85 (1961)．

\bibitem{kuwabara1972}
桑原尚夫，境久雄：連続音声中の切り出し母音および音節の音韻知覚，日本音響学会誌，Vol.28, No.5, pp.225--234 (1972)．

\bibitem{wichmann2001}
Wichmann, F. A. and Hill, N. J.: The psychometric function: I. Fitting, sampling, and goodness of fit, Perception \& Psychophysics, Vol.63, No.8, pp.1293--1313 (2001).

\bibitem{lakens2017}
Lakens, D.: Equivalence tests: A practical primer for t tests, correlations, and meta-analyses, Social Psychological and Personality Science, Vol.8, No.4, pp.355--362 (2017).

\bibitem{jomo_rules}
公益社団法人群馬県子ども会育成連合会：上毛かるた競技県大会規則（令和 7 年 8 月 21 日現在），\url{https://www.kodomo-kai.or.jp/gunma/}（参照 2026-09-24）．

\bibitem{voicevox}
ヒホ：VOICEVOX，\url{https://voicevox.hiroshiba.jp/}（参照 2026-10-08）．

\end{thebibliography}

\begin{biography}
\profile{m}{栗原 一貴}{津田塾大学学芸学部情報科学科教授．博士（情報理工学）．ヒューマンコンピュータインタラクション，エンタテインメントコンピューティングの研究に従事．}
% 丸山さんの所属と学年は投稿前に確認すること．
\profile{n}{丸山 礼華}{津田塾大学．競技かるたの聞き分けと動的フォントの研究に従事．}
\end{biography}

\end{document}
